\chapter{Friable approximants of \texorpdfstring{$\Xi$}{Xi} after Miyatake and the cosh ladder}\label{chap:miyatake}\label{ch:miyatake}\chaptermark{Friable approximants after Miyatake}

\ifdefined\conjecture\else\theoremstyle{sc}\newtheorem{conjecture}[theorem]{Conjecture}\theoremstyle{scr}\fi

Keep in Riemann's theta series only the integers composed of the primes up to $P$, and the result is a function $\Xi_P$ that differs from $\Xi$ by an exponentially small amount and yet has infinitely many zeros off the real line. This chapter studies that function, the friable approximant, as a test case for any computational criterion for the Riemann Hypothesis. It asks how many real zeros $\Xi_P$ shares with $\Xi$, why no finite set of primes reaches beyond a height of about $4q^2$, where $q$ is the first prime left out (a heuristic formula, checked by computation), where the remaining zeros go, and what the $2\sin^2$ lemma and the Thorin condition can and cannot detect. The last sections read an identity of Miyatake, who worked with the primes up to $P$ inside P\'olya's symmetrised kernel, probabilistically (L\'evy's area, a P\'olya Gamma clock, a logistic partition of unity), and place the bottom of the cosh ladder against the quasi Riemann abscissa.

Riemann's $\xi$ function is $\xi(s)=\tfrac12s(s-1)\pi^{-s/2}\Gamma(s/2)\zeta(s)$, and $\Xi(t)=\xi(\tfrac12+it)$ is even, entire and real on $\R$ (Chapter~\ref{chap:foundations}, Section~\ref{ch:clock}); RH says that all zeros of $\Xi$ are real. Miyatake \cite{Miyatake51b,Miyatake51a} approximated $\xi$ using only the integers composed of primes $p\le P$, so that the Dirichlet series becomes the finite Euler product $\zeta_P(s)=\prod_{p\le P}(1-p^{-s})^{-1}=\sum_{n\in\mathcal F(P)}n^{-s}$, with $\mathcal F(P)$ the set of $P$ friable integers. % UNCHECKED: Miyatake (the two 1951 papers could not be fetched; what they contain is stated here only as far as the title and the later sections of this chapter say)
He worked inside P\'olya's symmetrised kernel rather than Riemann's (Section~\ref{mi:sec:miyatake}). Here Riemann's own kernel is truncated instead, and the friable approximant $\Xi_P$ of Section~\ref{mi:sec:friable} is the object of study. The incomplete gamma form of truncated theta series, which is the form used for all computations below, is also the subject of Haglund \cite{Haglund09}. % UNCHECKED: Haglund (only the title of the paper was read; whether it treats the approximants studied here, weights w_n equal to 1 for n at most N, is to be settled by the authors)

The chapter is organised around one function, the deficit $D_P=\Xi_P-\Xi$, and one geometric object, the negative hump of $\Xi$. Every number printed below was computed by the scripts listed for this chapter in the code archive (file code/README.md), in floating point at the precision stated unless a certified method is named. The results are these.
\begin{enumerate}
\item $|\Xi_P-\Xi|\le|t^2+\tfrac14|\,\delta_P$ in $|\operatorname{Im}t|\le\tfrac32$ with $\delta_P\approx e^{-\pi q^2}/(\pi q^2)$, and $|\Xi_P-\Xi|\le(1+\delta(a_q))c_P$ on $\R$ with $c_P\approx4\pi q^2e^{-\pi q^2}$: the first omitted prime sets the accuracy (Theorems~\ref{mi:thm:approx} and~\ref{mi:thm:lorentz}, proved).
\item The truncated kernel has a kink of exact size $c_P>0$, so $\Xi_P\to c_P$ on $\R$; no reweighting of the theta series with $\sum_n|w_n|n^{-\sigma}<\infty$ for every $\sigma>0$ (this includes the friable sets) has infinitely many real zeros, and the same holds for any set of omitted integers (Theorems~\ref{mi:thm:kink} and~\ref{mi:thm:noproxy}, proved).
\item $D_P>0$ on $\R$ and is, with an explicit uniform error, a sum of Lorentzians (Theorems~\ref{mi:thm:positive} and~\ref{mi:thm:lorentz}, proved).
\item Real zeros of $\Xi_P$ live in negative humps of $\Xi$, two per surviving hump, and are lost in pairs by collision; the count is twice the number of humps deeper than $D_P$, at every height below $3\cdot10^{12}$ where the humps are short enough (Theorems~\ref{mi:thm:humps} and~\ref{mi:thm:count}, proved given the verification of RH to that height \cite{PT21}). For $P=7$ the computed count is $264$.
\item $\Xi_P$ has more zeros than $\Xi$: heuristically $N_P(X)=N(X)+y_P(X)/4+O(1)$, compared with argument principle counts for $P=2,3$, which fall below it by $4.86$ at $X=600$ for $P=2$ (Remark~\ref{mi:prop:zerocount}, heuristic and computed).
\item The $2\sin^2$ L\'evy density of $\Xi_2$ agrees with that of $\Xi$ to seven digits for $x\le0.1$ (computed), and on the computed range it is positive although $\Xi_2$ has zeros off the line; the support of the Thorin measure can (Section~\ref{mi:sec:sin2}).
\item Multiplying Riemann's kernel by $\cosh(bt)$ gives, unconditionally for $b\ge\frac12$ (Lagarias, Theorem~\ref{pol:thm:cosh}), reciprocals that are Laplace transforms of generalised gamma convolutions, and RH is equivalent to this persisting as $b\downarrow0$ (Section~\ref{mi:sec:ladder}).
\item Dividing instead by $\cosh(bt)$, as Miyatake does with $b=\frac92$, is smoothing by L\'evy's stochastic area, equivalently the de Bruijn Newman flow at a P\'olya Gamma random negative time; at $b=2$ the transform changes character, and for $b>2$ with $C_b\ne0$ it is never in the Laguerre P\'olya class (Sections~\ref{mi:sec:prob} and~\ref{mi:sec:btwo}).
\item The bottom of the cosh ladder is bounded by the quasi Riemann abscissa, and in paired form (real and imaginary parts of $\xi$ on vertical lines, a Thorin pair) it lies exactly at it (Section~\ref{mi:sec:quasi}).
\end{enumerate}
